\color{unidarkgreen}{\subsection{Максимальная токовая защита обмотки управления УШР (МТЗ ОУ)}}
\color{black}

\begin{enumerate}[label=\arabic{section}.\arabic{subsection}.\arabic*,  align=left, labelsep=4pt, leftmargin=0pt, itemindent=57pt]

\item
Защита предназначена для действия при КЗ на землю (корпус) обмотки управления УШР.
Функция использует измерения ТТ, установленного в цепи заземления (проводником) средней
точки обмотки управления.
\item
Функция может быть реализована в однофазном или трехфазном исполнении, данный выбор
задается программной накладкой «ИО МТЗ ОУ». Защита имеет регулируемые значения по току и
времени срабатывания. 

\begin{figure}[H]
\centering
\includegraphics[width=0.7\textwidth,height=0.7\textheight,keepaspectratio]{\fbpath/МТЗ ОУ/img/MTZ_OU.pdf}
\captionof{figure}{Функциональный блок МТЗ ОУ}\label{pict_mtzou:MTZ_OU}
\end{figure}

\small
\begin{longtable}{|>{\centering\arraybackslash}m{5.3cm}|>{\centering\arraybackslash}m{4.2cm}|>{\centering\arraybackslash}m{1cm}|>{\centering\arraybackslash}m{3.3cm}|>{\centering\arraybackslash}m{1.8cm}|}
% Первая страница: полный заголовок
\caption{Перечень уставочных параметров МТЗ ОУ\hfill\vspace{-0.5\baselineskip}}\label{mtz_ou:tbl1}\\
\hline
\rowcolor{gray!30}
Наименование уставки & Диапазон & Шаг & Значение по умолчанию & Единица измерения \\
\hline
\endfirsthead

% Промежуточные страницы: «Продолжение таблицы»
\caption*{\hspace{3pt}\emph{Продолжение таблицы \ref{mtz_ou:tbl1}}\hfill\ }\vspace{-0.5\baselineskip} \\
\hline
\rowcolor{gray!30}
Наименование уставки & Диапазон & Шаг & Значение по умолчанию & Единица измерения \\
\hline
\endhead

%===m>МТЗ ОУ|-
%===m
\end{longtable}
\normalsize

\end{enumerate}